Put
\[
S=\operatorname{span}\{r_0,r_1\}\subseteq\mathbb R^4.
\]
By \(\texttt{def:bow-target-coefficient}\), with the coordinates ordered by
\((a_{\mathrm b},z_{\mathrm b})\),
\[
r_0=\mathbf 1\otimes e_0,
\qquad
r_1=\mathbf 1\otimes e_1.
\]
Because \(e_0,e_1\) form the coordinate basis of \(\mathbb R^2\),
\[
S=\{\mathbf 1\otimes t:t\in\mathbb R^2\},
\qquad
\dim(S)=2.
\]
Indeed, the map \(t\mapsto\mathbf 1\otimes t\) is linear and injective,
since \(\mathbf 1\neq0\).

Let \(\mathcal I\) be an arbitrary menu satisfying the hypotheses.  The
inclusions \(r_0,r_1\in\operatorname{row}(A_{\mathrm b})\) imply
\[
S\subseteq\operatorname{row}(A_{\mathrm b}),
\qquad
\operatorname{rank}(A_{\mathrm b})\geq2.
\]
Suppose, toward a contradiction, that
\(\operatorname{rank}(A_{\mathrm b})<3\).  The preceding inequality and
\(\dim(S)=2\) then give
\[
\operatorname{rank}(A_{\mathrm b})=2,
\qquad
\operatorname{row}(A_{\mathrm b})=S.
\]
Fix \(\beta_X\in\mathcal L_{X_{\mathrm b}}\).  By
\(\texttt{def:bow-probe-matrix}\), reshaping the row indexed by \(\beta_X\)
as a matrix with output index \(a_{\mathrm b}\) and input index
\(z_{\mathrm b}\) yields
\[
R_{\beta_X}
=\phi_{X_{\mathrm b},\beta_X}
 \psi_{X_{\mathrm b},\beta_X}^{\mathsf T}.
\]
Every matrix represented by a vector in \(S\) has two identical rows.
Writing
\(\phi_{X_{\mathrm b},\beta_X}=(\phi(0),\phi(1))^{\mathsf T}\), the membership
\(R_{\beta_X}\in S\) therefore implies
\[
\bigl(\phi(0)-\phi(1)\bigr)
\psi_{X_{\mathrm b},\beta_X}^{\mathsf T}=0.
\]
The vector \(\psi_{X_{\mathrm b},\beta_X}\) is a probability state, so it is
nonnegative and its coordinates sum to one.  It is consequently nonzero,
and the last display forces \(\phi(0)=\phi(1)\).  Since \(\beta_X\) was
arbitrary, every treatment effect is a scalar multiple of \(\mathbf 1\),
whence
\[
\operatorname{span}\{\phi_{X_{\mathrm b},\beta_X}:
\beta_X\in\mathcal L_{X_{\mathrm b}}\}
\subseteq\operatorname{span}\{\mathbf 1\}
\neq\mathbb R^2.
\]
This contradicts the marginal effect-completeness hypothesis.  Hence
\[
\operatorname{rank}(A_{\mathrm b})\geq3.
\]
The rank of a matrix is at most its number of rows, so also
\[
\lvert\mathcal L_{X_{\mathrm b}}\rvert\geq3.
\]

It remains to prove attainment.  Under \(\texttt{def:binary-bow-menu}\), the
three treatment rows of \(A_{\mathrm b}^{\dagger}\) are
\[
w_0=\tfrac12e_0\otimes\psi^{+},
\qquad
w_1=\tfrac12e_1\otimes\psi^{+},
\qquad
w_2=\tfrac12\mathbf 1\otimes\psi^{-}.
\]
All three treatment effects are nonnegative and, using
\(\mathbf 1=e_0+e_1\),
\[
\tfrac12e_0+\tfrac12e_1+\tfrac12\mathbf 1
=e_0+e_1
=\mathbf 1.
\]
Thus the treatment instrument satisfies
\(\texttt{ass:instrument-normalization}\).  The terminal effects also obey
\(e_0+e_1=\mathbf 1\), so the terminal instrument is normalized as well.
Moreover,
\[
e_0=2\bigl(\tfrac12e_0\bigr),
\qquad
e_1=2\bigl(\tfrac12e_1\bigr),
\]
and therefore the treatment effects span \(\mathbb R^2\).

The two treatment states are linearly independent because
\[
\det\begin{pmatrix}
3/4&1/4\\
1/4&3/4
\end{pmatrix}
=\frac{9}{16}-\frac{1}{16}
=\frac12
\neq0.
\]
If \(\alpha w_0+\beta w_1+\gamma w_2=0\), comparison of the two
output-indexed blocks, followed by multiplication by two, gives
\[
\alpha\psi^{+}+\gamma\psi^{-}=0,
\qquad
\beta\psi^{+}+\gamma\psi^{-}=0.
\]
Linear independence of \(\psi^{+},\psi^{-}\) in the first equality yields
\(\alpha=\gamma=0\), and the second equality then yields \(\beta=0\).
Consequently the three rows are linearly independent and
\[
\operatorname{rank}(A_{\mathrm b}^{\dagger})=3.
\]

Finally, direct coordinate calculation gives
\[
e_0=\tfrac32\psi^{+}-\tfrac12\psi^{-},
\qquad
e_1=-\tfrac12\psi^{+}+\tfrac32\psi^{-}.
\]
Since
\[
w_0+w_1=\tfrac12\mathbf 1\otimes\psi^{+},
\qquad
w_2=\tfrac12\mathbf 1\otimes\psi^{-},
\]
we obtain the explicit row-space representations
\[
r_0=\mathbf 1\otimes e_0=3(w_0+w_1)-w_2,
\qquad
r_1=\mathbf 1\otimes e_1=-(w_0+w_1)+3w_2.
\]
Thus \(r_0,r_1\in\operatorname{row}(A_{\mathrm b}^{\dagger})\).  The
universal lower bound and this normalized, marginally effect-complete
rank-three witness prove that the minimum attainable treatment branch-matrix
rank is exactly three.  The proof uses the target-row assumption in the
lower-bound argument and therefore makes no claim that every rank-three menu
contains \(r_0\) and \(r_1\).